%============================================================
%  Bd. 22 - Gerichtete Graphen
%============================================================

\documentclass[main.tex]{subfiles}

\ifSubfilesClassLoaded{
  \BandSetupStandalone{B22}
}{
}

\title{Bd. 22 - Gerichtete Graphen}
\author{Martin Kunze}
\date{}
\setcounter{file}{22}

\hypersetup{
  pdftitle={Bd. 22 - Gerichtete Graphen},
  pdfauthor={Martin Kunze}
}

\begin{document}

\title{Bd. 22 - Gerichtete Graphen}
\author{Martin Kunze}
\date{}
\hypersetup{pageanchor=false}
\maketitle
\hypersetup{pageanchor=true}
\setcounter{tocdepth}{3}
\tableofcontents
\setcounter{file}{22}
\renewcommand{\FormulaBandID}{22}

\chapter{Einleitung}

Ein gerichteter Graph besteht aus einer Knotenmenge und einer Relation auf
dieser Menge. Die mengentheoretische Grunddefinition steht bereits in
Band~03; der vorliegende Band entwickelt daraus die elementaren Begriffe der
Graphentheorie. Ein geordnetes Paar \((x,y)\) in der Kantenrelation wird als
von \(x\) nach \(y\) gerichteter Bogen gelesen. Schleifen sind zunächst
zugelassen. Mehrfachbögen treten in der relationalen Darstellung nicht auf.

Endliche Kantenzüge werden als nullbasierte endliche Folgen im Sinne von
Band~21 erfasst. Dadurch bleiben Knotenfolge, Kantenlänge und Endpunkte
explizit typisiert. Band~23 ergänzt die Symmetrieaxiome ungerichteter Graphen,
Band~24 deren schlichte und Band~25 deren zusammenhängende Spezialisierung.
Band~26 führt darauf aufbauend Bäume ein.

\chapter{Grundbegriffe gerichteter Graphen}

\section{Graphen und Endpunkte}

Wir schreiben \(\GraphStruct{V,E}\), wenn \(E\) eine Kantenrelation auf
dem Knotenträger \(V\) ist. Nach der Definition aus Band~03 bedeutet dies
genau \(E\subseteq V\times V\).

\FormulaThmDeltaK[Zugriff auf die Kantenrelation eines Graphen]{%
  \GraphStruct{V,E}\vdash E\subseteq V\times V%
}{GraphEdgeRelationAccess}{%
  \DeltaRow{Mengen}{V\dsep E}%
}
\begin{tabproof}
  \proofstep{1}{\GraphStruct{V,E}}{\rA}
  \proofstep{1}{E\subseteq V\times V}{%
    \FormulaRefAuto{GraphDef}[def]{1}}
\end{tabproof}

\FormulaThmDeltaK[Typisierung der Endpunkte eines Bogens]{%
  \GraphStruct{V,E}\dsep (x,y)\in E\vdash
  x\in V\land y\in V%
}{GraphEdgeEndpointTyping}{%
  \DeltaRow{Mengen}{V\dsep E\dsep x\dsep y}%
}
\begin{tabproof}
  \proofstep{1}{\GraphStruct{V,E}}{\rA}
  \proofstep{2}{(x,y)\in E}{\rA}
  \proofstep{1}{E\subseteq V\times V}{%
    \FormulaRefAuto{GraphEdgeRelationAccess}[thm]{1}}
  \proofstep{1,2}{(x,y)\in V\times V}{%
    \FormulaRefAuto{A\subseteq B\dsep x\in A\vdash x\in B}[thm]{3,2}}
  \proofstep{1,2}{x\in V\land y\in V}{%
    \FormulaRefAuto{(a,b)\in A\times B\eqvdash
      a\in A\land b\in B}[thm]{4}}
\end{tabproof}

\section{Ein- und Ausnachbarschaften}

\FormulaDefDeltaK[Ausnachbarschaft eines Knotens]{%
  \GraphStruct{V,E}\dsep x\in V\vdash
  \OutNeighbors{E}{x}\coloneqq
  \{y\in V\mid (x,y)\in E\}%
}{DirectedOutNeighborhoodDef}{%
  \DeltaRow{Mengen}{V\dsep E\dsep x\dsep y\dsep\OutNeighbors{E}{x}}
  \DeltaRow{\textbf{Neues Symbol}}{}
  \DeltaPrem{Ausnachbarschaft}{\OutNeighbors{E}{x}}
}

\FormulaDefDeltaK[Einnachbarschaft eines Knotens]{%
  \GraphStruct{V,E}\dsep x\in V\vdash
  \InNeighbors{E}{x}\coloneqq
  \{y\in V\mid (y,x)\in E\}%
}{DirectedInNeighborhoodDef}{%
  \DeltaRow{Mengen}{V\dsep E\dsep x\dsep y\dsep\InNeighbors{E}{x}}
  \DeltaRow{\textbf{Neues Symbol}}{}
  \DeltaPrem{Einnachbarschaft}{\InNeighbors{E}{x}}
}

\FormulaThmDeltaK[Elementkriterium der Ausnachbarschaft]{%
  \GraphStruct{V,E}\dsep x,y\in V\vdash
  y\in\OutNeighbors{E}{x}\leftrightarrow (x,y)\in E%
}{DirectedOutNeighborhoodMembership}{%
  \DeltaRow{Mengen}{V\dsep E\dsep x\dsep y}%
}
\begin{tabproofsplitwide}
  \proofpartwideindR[\(\rightarrow\)]{%
    \GraphStruct{V,E}\dsep x,y\in V\vdash
    y\in\OutNeighbors{E}{x}\rightarrow (x,y)\in E
  }{%
    \GraphStruct{V,E}\dsep x,y\in V\vdash
    y\in\OutNeighbors{E}{x}\rightarrow (x,y)\in E
  }[DirectedOutNeighborhoodMembershipForward]
    \proofstepwidestar[1]{\GraphStruct{V,E}}{\rA}
    \proofstepwidestar[2]{x\in V}{\rA}
    \proofstepwidestar[3]{y\in V}{\rA}
    \proofstepwidestar[4]{y\in\OutNeighbors{E}{x}}{\rA}
    \proofstepwidestar[1,2]{%
      \OutNeighbors{E}{x}=\{z\in V\mid(x,z)\in E\}}{%
      \FormulaRefAuto{DirectedOutNeighborhoodDef}[def]{1,2}}
    \proofstepwidestar[1,2,4]{y\in\{z\in V\mid(x,z)\in E\}}{%
      \rIE{5,4}}
    \proofstepwidestar[1,2,4]{(x,y)\in E}{%
      \FormulaRefAuto{x\in\{x\in A\mid P(x)\}\vdash P(x)}[thm]{6}}
    \proofstepwidestar[1,2,3]{%
      y\in\OutNeighbors{E}{x}\rightarrow(x,y)\in E}{\rRI{4,7}}
  \closeproofpartwideindR

  \proofpartwideindR[\(\leftarrow\)]{%
    \GraphStruct{V,E}\dsep x,y\in V\vdash
    (x,y)\in E\rightarrow y\in\OutNeighbors{E}{x}
  }{%
    \GraphStruct{V,E}\dsep x,y\in V\vdash
    (x,y)\in E\rightarrow y\in\OutNeighbors{E}{x}
  }[DirectedOutNeighborhoodMembershipBackward]
    \proofstepwidestar[1]{\GraphStruct{V,E}}{\rA}
    \proofstepwidestar[2]{x\in V}{\rA}
    \proofstepwidestar[3]{y\in V}{\rA}
    \proofstepwidestar[4]{(x,y)\in E}{\rA}
    \proofstepwidestar[1,2]{%
      \OutNeighbors{E}{x}=\{z\in V\mid(x,z)\in E\}}{%
      \FormulaRefAuto{DirectedOutNeighborhoodDef}[def]{1,2}}
    \proofstepwidestar[3,4]{y\in\{z\in V\mid(x,z)\in E\}}{%
      \FormulaRefAuto{x\in A\dsep P(x)\vdash
        x\in\{x\in A\mid P(x)\}}[thm]{3,4}}
    \proofstepwidestar[1,2]{%
      \{z\in V\mid(x,z)\in E\}=\OutNeighbors{E}{x}}{%
      \FormulaRefAuto{a=b\vdash b=a}[thm]{5}}
    \proofstepwidestar[1,2,3,4]{y\in\OutNeighbors{E}{x}}{\rIE{7,6}}
    \proofstepwidestar[1,2,3]{%
      (x,y)\in E\rightarrow y\in\OutNeighbors{E}{x}}{\rRI{4,8}}
  \closeproofpartwideindR

  \proofpartwide{Schluss}
    \proofstepwidestar[1]{\GraphStruct{V,E}}{\rA}
    \proofstepwidestar[2]{x\in V}{\rA}
    \proofstepwidestar[3]{y\in V}{\rA}
    \proofstepwidestar[1,2,3]{%
      y\in\OutNeighbors{E}{x}\rightarrow(x,y)\in E}{%
      \FormulaRefAuto{DirectedOutNeighborhoodMembershipForward}[thm]{1,2,3}}
    \proofstepwidestar[1,2,3]{%
      (x,y)\in E\rightarrow y\in\OutNeighbors{E}{x}}{%
      \FormulaRefAuto{DirectedOutNeighborhoodMembershipBackward}[thm]{1,2,3}}
    \proofstepwidestar[1,2,3]{%
      y\in\OutNeighbors{E}{x}\leftrightarrow(x,y)\in E}{\rLRI{4,5}}
  \closeproofpartwide
\end{tabproofsplitwide}

\FormulaThmDeltaK[Elementkriterium der Einnachbarschaft]{%
  \GraphStruct{V,E}\dsep x,y\in V\vdash
  y\in\InNeighbors{E}{x}\leftrightarrow (y,x)\in E%
}{DirectedInNeighborhoodMembership}{%
  \DeltaRow{Mengen}{V\dsep E\dsep x\dsep y}%
}
\begin{tabproofsplitwide}
  \proofpartwideindR[\(\rightarrow\)]{%
    \GraphStruct{V,E}\dsep x,y\in V\vdash
    y\in\InNeighbors{E}{x}\rightarrow(y,x)\in E
  }{%
    \GraphStruct{V,E}\dsep x,y\in V\vdash
    y\in\InNeighbors{E}{x}\rightarrow(y,x)\in E
  }[DirectedInNeighborhoodMembershipForward]
    \proofstepwidestar[1]{\GraphStruct{V,E}}{\rA}
    \proofstepwidestar[2]{x\in V}{\rA}
    \proofstepwidestar[3]{y\in V}{\rA}
    \proofstepwidestar[4]{y\in\InNeighbors{E}{x}}{\rA}
    \proofstepwidestar[1,2]{%
      \InNeighbors{E}{x}=\{z\in V\mid(z,x)\in E\}}{%
      \FormulaRefAuto{DirectedInNeighborhoodDef}[def]{1,2}}
    \proofstepwidestar[1,2,4]{y\in\{z\in V\mid(z,x)\in E\}}{\rIE{5,4}}
    \proofstepwidestar[1,2,4]{(y,x)\in E}{%
      \FormulaRefAuto{x\in\{x\in A\mid P(x)\}\vdash P(x)}[thm]{6}}
    \proofstepwidestar[1,2,3]{%
      y\in\InNeighbors{E}{x}\rightarrow(y,x)\in E}{\rRI{4,7}}
  \closeproofpartwideindR

  \proofpartwideindR[\(\leftarrow\)]{%
    \GraphStruct{V,E}\dsep x,y\in V\vdash
    (y,x)\in E\rightarrow y\in\InNeighbors{E}{x}
  }{%
    \GraphStruct{V,E}\dsep x,y\in V\vdash
    (y,x)\in E\rightarrow y\in\InNeighbors{E}{x}
  }[DirectedInNeighborhoodMembershipBackward]
    \proofstepwidestar[1]{\GraphStruct{V,E}}{\rA}
    \proofstepwidestar[2]{x\in V}{\rA}
    \proofstepwidestar[3]{y\in V}{\rA}
    \proofstepwidestar[4]{(y,x)\in E}{\rA}
    \proofstepwidestar[1,2]{%
      \InNeighbors{E}{x}=\{z\in V\mid(z,x)\in E\}}{%
      \FormulaRefAuto{DirectedInNeighborhoodDef}[def]{1,2}}
    \proofstepwidestar[3,4]{y\in\{z\in V\mid(z,x)\in E\}}{%
      \FormulaRefAuto{x\in A\dsep P(x)\vdash
        x\in\{x\in A\mid P(x)\}}[thm]{3,4}}
    \proofstepwidestar[1,2]{%
      \{z\in V\mid(z,x)\in E\}=\InNeighbors{E}{x}}{%
      \FormulaRefAuto{a=b\vdash b=a}[thm]{5}}
    \proofstepwidestar[1,2,3,4]{y\in\InNeighbors{E}{x}}{\rIE{7,6}}
    \proofstepwidestar[1,2,3]{%
      (y,x)\in E\rightarrow y\in\InNeighbors{E}{x}}{\rRI{4,8}}
  \closeproofpartwideindR

  \proofpartwide{Schluss}
    \proofstepwidestar[1]{\GraphStruct{V,E}}{\rA}
    \proofstepwidestar[2]{x\in V}{\rA}
    \proofstepwidestar[3]{y\in V}{\rA}
    \proofstepwidestar[1,2,3]{%
      y\in\InNeighbors{E}{x}\rightarrow(y,x)\in E}{%
      \FormulaRefAuto{DirectedInNeighborhoodMembershipForward}[thm]{1,2,3}}
    \proofstepwidestar[1,2,3]{%
      (y,x)\in E\rightarrow y\in\InNeighbors{E}{x}}{%
      \FormulaRefAuto{DirectedInNeighborhoodMembershipBackward}[thm]{1,2,3}}
    \proofstepwidestar[1,2,3]{%
      y\in\InNeighbors{E}{x}\leftrightarrow(y,x)\in E}{\rLRI{4,5}}
  \closeproofpartwide
\end{tabproofsplitwide}

\FormulaThmDeltaK[Gebündelte Nachbarschaftskriterien]{%
  \GraphStruct{V,E}\dsep x,y\in V\vdash
  \begin{aligned}[t]
    &\bigl(y\in\OutNeighbors{E}{x}\leftrightarrow(x,y)\in E\bigr)\land{}\\[-2pt]
    &\bigl(y\in\InNeighbors{E}{x}\leftrightarrow(y,x)\in E\bigr)
  \end{aligned}%
}{DirectedNeighborhoodMembership}{%
  \DeltaRow{Mengen}{V\dsep E\dsep x\dsep y}%
}
\begin{tabproofwide}
  \proofstepwidestar[1]{\GraphStruct{V,E}}{\rA}
  \proofstepwidestar[2]{x\in V}{\rA}
  \proofstepwidestar[3]{y\in V}{\rA}
  \proofstepwidestar[1,2,3]{%
    y\in\OutNeighbors{E}{x}\leftrightarrow(x,y)\in E}{%
    \FormulaRefAuto{DirectedOutNeighborhoodMembership}[thm]{1,2,3}}
  \proofstepwidestar[1,2,3]{%
    y\in\InNeighbors{E}{x}\leftrightarrow(y,x)\in E}{%
    \FormulaRefAuto{DirectedInNeighborhoodMembership}[thm]{1,2,3}}
  \proofstepwidestar[1,2,3]{%
    \bigl(y\in\OutNeighbors{E}{x}\leftrightarrow(x,y)\in E\bigr)\land
    \bigl(y\in\InNeighbors{E}{x}\leftrightarrow(y,x)\in E\bigr)}{\rAI{4,5}}
\end{tabproofwide}

\section{Umkehrgraph und induzierte Teilgraphen}

\FormulaDefDeltaK[Umkehrrelation eines Graphen]{%
  \GraphStruct{V,E}\vdash
  \ReverseEdges{E}\coloneqq
  \{(x,y)\in V\times V\mid(y,x)\in E\}%
}{DirectedReverseEdgeRelationDef}{%
  \DeltaRow{Mengen}{V\dsep E\dsep x\dsep y\dsep\ReverseEdges{E}}
  \DeltaRow{\textbf{Neues Symbol}}{}
  \DeltaPrem{Umkehrrelation}{\ReverseEdges{E}}
}

\FormulaThmDeltaK[Der Umkehrgraph ist ein Graph]{%
  \GraphStruct{V,E}\vdash\GraphStruct{V,\ReverseEdges{E}}%
}{DirectedReverseGraph}{%
  \DeltaRow{Mengen}{V\dsep E\dsep\ReverseEdges{E}}%
}
\begin{tabproof}
  \proofstep{1}{\GraphStruct{V,E}}{\rA}
  \proofstep{1}{%
    \ReverseEdges{E}=\{(x,y)\in V\times V\mid(y,x)\in E\}}{%
    \FormulaRefAuto{DirectedReverseEdgeRelationDef}[def]{1}}
  \proofstep{}{\{(x,y)\in V\times V\mid(y,x)\in E\}\subseteq V\times V}{%
    \FormulaRefAuto{\{x\in A\mid P(x)\}\subseteq A}[thm]}
  \proofstep{1}{\ReverseEdges{E}\subseteq V\times V}{\rIE{2,3}}
  \proofstep{1}{\GraphStruct{V,\ReverseEdges{E}}}{%
    \FormulaRefAuto{GraphDef}[def]{4}}
\end{tabproof}

\FormulaThmDeltaK[Elementkriterium des Umkehrgraphen]{%
  \GraphStruct{V,E}\dsep x,y\in V\vdash
  (x,y)\in\ReverseEdges{E}\leftrightarrow(y,x)\in E%
}{DirectedReverseEdgeMembership}{%
  \DeltaRow{Mengen}{V\dsep E\dsep x\dsep y}%
}
\begin{tabproofsplitwide}
  \proofpartwideindR[\(\rightarrow\)]{%
    \GraphStruct{V,E}\dsep x,y\in V\vdash
    (x,y)\in\ReverseEdges{E}\rightarrow(y,x)\in E
  }{%
    \GraphStruct{V,E}\dsep x,y\in V\vdash
    (x,y)\in\ReverseEdges{E}\rightarrow(y,x)\in E
  }[DirectedReverseEdgeMembershipForward]
    \proofstepwidestar[1]{\GraphStruct{V,E}}{\rA}
    \proofstepwidestar[2]{x\in V}{\rA}
    \proofstepwidestar[3]{y\in V}{\rA}
    \proofstepwidestar[4]{(x,y)\in\ReverseEdges{E}}{\rA}
    \proofstepwidestar[1]{%
      \ReverseEdges{E}=\{(u,v)\in V\times V\mid(v,u)\in E\}}{%
      \FormulaRefAuto{DirectedReverseEdgeRelationDef}[def]{1}}
    \proofstepwidestar[1,4]{%
      (x,y)\in\{(u,v)\in V\times V\mid(v,u)\in E\}}{\rIE{5,4}}
    \proofstepwidestar[1,4]{(y,x)\in E}{%
      \FormulaRefAuto{x\in\{x\in A\mid P(x)\}\vdash P(x)}[thm]{6}}
    \proofstepwidestar[1,2,3]{%
      (x,y)\in\ReverseEdges{E}\rightarrow(y,x)\in E}{\rRI{4,7}}
  \closeproofpartwideindR

  \proofpartwideindR[\(\leftarrow\)]{%
    \GraphStruct{V,E}\dsep x,y\in V\vdash
    (y,x)\in E\rightarrow(x,y)\in\ReverseEdges{E}
  }{%
    \GraphStruct{V,E}\dsep x,y\in V\vdash
    (y,x)\in E\rightarrow(x,y)\in\ReverseEdges{E}
  }[DirectedReverseEdgeMembershipBackward]
    \proofstepwidestar[1]{\GraphStruct{V,E}}{\rA}
    \proofstepwidestar[2]{x\in V}{\rA}
    \proofstepwidestar[3]{y\in V}{\rA}
    \proofstepwidestar[4]{(y,x)\in E}{\rA}
    \proofstepwidestar[2,3]{(x,y)\in V\times V}{%
      \FormulaRefAuto{a\in A\dsep b\in B\vdash(a,b)\in A\times B}[thm]{2,3}}
    \proofstepwidestar[2,3,4]{%
      (x,y)\in\{(u,v)\in V\times V\mid(v,u)\in E\}}{%
      \FormulaRefAuto{x\in A\dsep P(x)\vdash
        x\in\{x\in A\mid P(x)\}}[thm]{5,4}}
    \proofstepwidestar[1]{%
      \{(u,v)\in V\times V\mid(v,u)\in E\}=\ReverseEdges{E}}{%
      \FormulaRefAuto{a=b\vdash b=a}[thm]{%
        \FormulaRefAuto{DirectedReverseEdgeRelationDef}[def]{1}}}
    \proofstepwidestar[1,2,3,4]{(x,y)\in\ReverseEdges{E}}{\rIE{7,6}}
    \proofstepwidestar[1,2,3]{%
      (y,x)\in E\rightarrow(x,y)\in\ReverseEdges{E}}{\rRI{4,8}}
  \closeproofpartwideindR

  \proofpartwide{Schluss}
    \proofstepwidestar[1]{\GraphStruct{V,E}}{\rA}
    \proofstepwidestar[2]{x\in V}{\rA}
    \proofstepwidestar[3]{y\in V}{\rA}
    \proofstepwidestar[1,2,3]{%
      (x,y)\in\ReverseEdges{E}\rightarrow(y,x)\in E}{%
      \FormulaRefAuto{DirectedReverseEdgeMembershipForward}[thm]{1,2,3}}
    \proofstepwidestar[1,2,3]{%
      (y,x)\in E\rightarrow(x,y)\in\ReverseEdges{E}}{%
      \FormulaRefAuto{DirectedReverseEdgeMembershipBackward}[thm]{1,2,3}}
    \proofstepwidestar[1,2,3]{%
      (x,y)\in\ReverseEdges{E}\leftrightarrow(y,x)\in E}{\rLRI{4,5}}
  \closeproofpartwide
\end{tabproofsplitwide}

\FormulaDefDeltaK[Induzierte Kantenrelation]{%
  \GraphStruct{V,E}\dsep U\subseteq V\vdash
  \DirectedInducedEdges{E}{U}\coloneqq E\cap(U\times U)%
}{DirectedInducedEdgeRelationDef}{%
  \DeltaRow{Mengen}{V\dsep E\dsep U\dsep\DirectedInducedEdges{E}{U}}
  \DeltaRow{\textbf{Neues Symbol}}{}
  \DeltaPrem{Induzierte Kantenrelation}{\DirectedInducedEdges{E}{U}}
}

\FormulaThmDeltaK[Der induzierte Teilgraph ist ein Graph]{%
  \GraphStruct{V,E}\dsep U\subseteq V\vdash
  \GraphStruct{U,\DirectedInducedEdges{E}{U}}%
}{DirectedInducedSubgraph}{%
  \DeltaRow{Mengen}{V\dsep E\dsep U}%
}
\begin{tabproof}
  \proofstep{1}{\GraphStruct{V,E}}{\rA}
  \proofstep{2}{U\subseteq V}{\rA}
  \proofstep{1,2}{\DirectedInducedEdges{E}{U}=E\cap(U\times U)}{%
    \FormulaRefAuto{DirectedInducedEdgeRelationDef}[def]{1,2}}
  \proofstep{}{E\cap(U\times U)\subseteq U\times U}{%
    \FormulaRefAuto{A\cap B\subseteq B}[thm]}
  \proofstep{1,2}{\DirectedInducedEdges{E}{U}\subseteq U\times U}{\rIE{3,4}}
  \proofstep{1,2}{\GraphStruct{U,\DirectedInducedEdges{E}{U}}}{%
    \FormulaRefAuto{GraphDef}[def]{5}}
\end{tabproof}

\FormulaThmDeltaK[Elementkriterium der induzierten Kantenrelation]{%
  \GraphStruct{V,E}\dsep U\subseteq V\vdash
  \begin{aligned}[t]
    &(x,y)\in\DirectedInducedEdges{E}{U}\leftrightarrow{}\\[-2pt]
    &\qquad (x,y)\in E\land x\in U\land y\in U
  \end{aligned}%
}{DirectedInducedEdgeMembership}{%
  \DeltaRow{Mengen}{V\dsep E\dsep U\dsep x\dsep y}%
}
\begin{tabproofsplitwide}
  \proofpartwideindR[\(\rightarrow\)]{%
    \begin{aligned}[t]
    &\GraphStruct{V,E}\dsep U\subseteq V\vdash{}\\[-2pt]
    &\quad (x,y)\in\DirectedInducedEdges{E}{U}\rightarrow
      \bigl((x,y)\in E\land x\in U\land y\in U\bigr)
    \end{aligned}
  }{%
    \GraphStruct{V,E}\dsep U\subseteq V\vdash
    (x,y)\in\DirectedInducedEdges{E}{U}\rightarrow
    \bigl((x,y)\in E\land x\in U\land y\in U\bigr)
  }[DirectedInducedEdgeMembershipForward]
    \proofstepwidestar[1]{\GraphStruct{V,E}}{\rA}
    \proofstepwidestar[2]{U\subseteq V}{\rA}
    \proofstepwidestar[3]{(x,y)\in\DirectedInducedEdges{E}{U}}{\rA}
    \proofstepwidestar[1,2]{\DirectedInducedEdges{E}{U}=E\cap(U\times U)}{%
      \FormulaRefAuto{DirectedInducedEdgeRelationDef}[def]{1,2}}
    \proofstepwidestar[1,2,3]{(x,y)\in E\cap(U\times U)}{\rIE{4,3}}
    \proofstepwidestar[1,2,3]{(x,y)\in E\land(x,y)\in U\times U}{%
      \FormulaRefAuto{x\in A\cap B\eqvdash x\in A\land x\in B}[thm]{5}}
    \proofstepwidestar[1,2,3]{(x,y)\in E}{\rAEa{6}}
    \proofstepwidestar[1,2,3]{(x,y)\in U\times U}{\rAEb{6}}
    \proofstepwidestar[1,2,3]{x\in U\land y\in U}{%
      \FormulaRefAuto{(a,b)\in A\times B\eqvdash
        a\in A\land b\in B}[thm]{8}}
    \proofstepwidestar[1,2,3]{(x,y)\in E\land x\in U\land y\in U}{%
      \FormulaRefAuto{P\dsep Q\dsep R\vdash P\land Q\land R}[thm]{%
        7,\rAEa{9},\rAEb{9}}}
    \proofstepwidestarbreak[1,2]{%
      (x,y)\in\DirectedInducedEdges{E}{U}\rightarrow
      \bigl((x,y)\in E\land x\in U\land y\in U\bigr)}{\rRI{3,10}}
  \closeproofpartwideindR

  \proofpartwideindR[\(\leftarrow\)]{%
    \begin{aligned}[t]
    &\GraphStruct{V,E}\dsep U\subseteq V\vdash{}\\[-2pt]
    &\quad\bigl((x,y)\in E\land x\in U\land y\in U\bigr)
      \rightarrow(x,y)\in\DirectedInducedEdges{E}{U}
    \end{aligned}
  }{%
    \GraphStruct{V,E}\dsep U\subseteq V\vdash
    \bigl((x,y)\in E\land x\in U\land y\in U\bigr)
    \rightarrow(x,y)\in\DirectedInducedEdges{E}{U}
  }[DirectedInducedEdgeMembershipBackward]
    \proofstepwidestar[1]{\GraphStruct{V,E}}{\rA}
    \proofstepwidestar[2]{U\subseteq V}{\rA}
    \proofstepwidestar[3]{(x,y)\in E\land x\in U\land y\in U}{\rA}
    \proofstepwidestar[3]{(x,y)\in E}{\rAEa{3}}
    \proofstepwidestar[3]{x\in U}{%
      \FormulaRefAuto{P\land Q\land R\vdash Q}[thm]{3}}
    \proofstepwidestar[3]{y\in U}{%
      \FormulaRefAuto{P\land Q\land R\vdash R}[thm]{3}}
    \proofstepwidestar[3]{(x,y)\in U\times U}{%
      \FormulaRefAuto{a\in A\dsep b\in B\vdash(a,b)\in A\times B}[thm]{5,6}}
    \proofstepwidestar[3]{(x,y)\in E\cap(U\times U)}{%
      \FormulaRefAuto{x\in A\dsep x\in B\vdash x\in A\cap B}[thm]{4,7}}
    \proofstepwidestar[1,2]{E\cap(U\times U)=\DirectedInducedEdges{E}{U}}{%
      \FormulaRefAuto{a=b\vdash b=a}[thm]{%
        \FormulaRefAuto{DirectedInducedEdgeRelationDef}[def]{1,2}}}
    \proofstepwidestar[1,2,3]{(x,y)\in\DirectedInducedEdges{E}{U}}{\rIE{9,8}}
    \proofstepwidestarbreak[1,2]{%
      \bigl((x,y)\in E\land x\in U\land y\in U\bigr)
      \rightarrow(x,y)\in\DirectedInducedEdges{E}{U}}{\rRI{3,10}}
  \closeproofpartwideindR

  \proofpartwide{Schluss}
    \proofstepwidestar[1]{\GraphStruct{V,E}}{\rA}
    \proofstepwidestar[2]{U\subseteq V}{\rA}
    \proofstepwidestar[1,2]{%
      (x,y)\in\DirectedInducedEdges{E}{U}\rightarrow
      \bigl((x,y)\in E\land x\in U\land y\in U\bigr)}{%
      \FormulaRefAuto{DirectedInducedEdgeMembershipForward}[thm]{1,2}}
    \proofstepwidestar[1,2]{%
      \bigl((x,y)\in E\land x\in U\land y\in U\bigr)
      \rightarrow(x,y)\in\DirectedInducedEdges{E}{U}}{%
      \FormulaRefAuto{DirectedInducedEdgeMembershipBackward}[thm]{1,2}}
    \proofstepwidestar[1,2]{%
      (x,y)\in\DirectedInducedEdges{E}{U}\leftrightarrow
      \bigl((x,y)\in E\land x\in U\land y\in U\bigr)}{\rLRI{3,4}}
  \closeproofpartwide
\end{tabproofsplitwide}

\chapter{Endliche Kantenzüge und Erreichbarkeit}

\section{Walks und Pfade}

\FormulaDefDeltaK[Endlicher gerichteter Walk]{%
  \DirectedWalk{V}{E}{p}{n}%
}{DirectedWalkDef}{%
  \DeltaRow{Mengen}{V\dsep E\dsep p\dsep n\dsep i}
  \DeltaPrem{Graph}{\GraphStruct{V,E}}
  \DeltaRow{\textbf{Axiome}}{}
  \DeltaRow{Natürliche Kantenlänge}{%
    n\in\mathbb N}
  \DeltaRow{Knotenfolge}{%
    p\colon\NatSeg{n+1}\to V}
  \DeltaRow{Kantenschritt}{%
    \forall i\in\NatSeg n\,((p(i),p(i+1))\in E)}
  \DeltaRow{\textbf{Neues Symbol}}{}
  \DeltaPrem{Struktur}{\DirectedWalk{V}{E}{p}{n}}
}

\FormulaAxiomDeltaK[Natürliche Kantenlänge]{%
  \begin{aligned}[t]
    &\GraphStruct{V,E}\dsep \DirectedWalk{V}{E}{p}{n}\vdash{}\\[-2pt]
    &n\in\mathbb N
  \end{aligned}%
}{DirectedWalkLengthAxiom}{%
  \DeltaRow{Mengen}{V\dsep E\dsep p\dsep n\dsep i}
}

\FormulaAxiomDeltaK[Knotenfolge]{%
  \begin{aligned}[t]
    &\GraphStruct{V,E}\dsep \DirectedWalk{V}{E}{p}{n}\vdash{}\\[-2pt]
    &p\colon\NatSeg{n+1}\to V
  \end{aligned}%
}{DirectedWalkSequenceAxiom}{%
  \DeltaRow{Mengen}{V\dsep E\dsep p\dsep n\dsep i}
}

\FormulaAxiomDeltaK[Aufeinanderfolgende Kanten]{%
  \begin{aligned}[t]
    &\GraphStruct{V,E}\dsep \DirectedWalk{V}{E}{p}{n}\vdash{}\\[-2pt]
    &\forall i\in\NatSeg n\,((p(i),p(i+1))\in E)
  \end{aligned}%
}{DirectedWalkEdgesAxiom}{%
  \DeltaRow{Mengen}{V\dsep E\dsep p\dsep n\dsep i}
}

\FormulaThmDeltaK[Charakterisierung durch die Axiome]{%
  \GraphStruct{V,E}\vdash
  \begin{aligned}[t]
    \DirectedWalk{V}{E}{p}{n}\eqvdash{}&
      n\in\mathbb N\land
      p\colon\NatSeg{n+1}\to V\land{}\\[-2pt]
      &\forall i\in\NatSeg n\,
        \bigl((p(i),p(i+1))\in E\bigr)
  \end{aligned}%
}{DirectedWalkCharacterization}{%
  \DeltaRow{Mengen}{V\dsep E\dsep p\dsep n\dsep i}
}
\begin{tabproofsplitwide}
  \proofpartwide{Hinrichtung}
    \proofstepwidestar[1]{\GraphStruct{V,E}}{\rA}
    \proofstepwidestar[2]{\DirectedWalk{V}{E}{p}{n}}{\rA}
    \proofstepwidestar[1,2]{n\in\mathbb N}{%
      \FormulaRefAuto{DirectedWalkLengthAxiom}[ax]{1,2}}
    \proofstepwidestar[1,2]{p\colon\NatSeg{n+1}\to V}{%
      \FormulaRefAuto{DirectedWalkSequenceAxiom}[ax]{1,2}}
    \proofstepwidestar[1,2]{\forall i\in\NatSeg n\,((p(i),p(i+1))\in E)}{%
      \FormulaRefAuto{DirectedWalkEdgesAxiom}[ax]{1,2}}
    \proofstepwidestar[1,2]{%
      \begin{aligned}[t]
      &n\in\mathbb N\land{}\\[-2pt]
      &p\colon\NatSeg{n+1}\to V\land{}\\[-2pt]
      &\forall i\in\NatSeg n\,((p(i),p(i+1))\in E)
    \end{aligned}}{\rAIStar{3,4,5}}
  \closeproofpartwide

  \proofpartwide{Rückrichtung}
    \proofstepwidestar[1]{\GraphStruct{V,E}}{\rA}
    \proofstepwidestar[2]{%
      \begin{aligned}[t]
      &n\in\mathbb N\land{}\\[-2pt]
      &p\colon\NatSeg{n+1}\to V\land{}\\[-2pt]
      &\forall i\in\NatSeg n\,((p(i),p(i+1))\in E)
    \end{aligned}}{\rA}
    \proofstepwidestar[2]{n\in\mathbb N}{\rAEa{2}}
    \proofstepwidestar[2]{p\colon\NatSeg{n+1}\to V}{\rAEa{\rAEb{2}}}
    \proofstepwidestar[2]{\forall i\in\NatSeg n\,((p(i),p(i+1))\in E)}{\rAEb{\rAEb{2}}}
    \proofstepwidestar[1,2]{\DirectedWalk{V}{E}{p}{n}}{%
      \FormulaRefAuto{DirectedWalkDef}[def]{1,3,4,5}}
  \closeproofpartwide
\end{tabproofsplitwide}

Die Zahl \(n\) ist die Kantenlänge des Walks; seine Knotenfolge besitzt
den Bereich \(\NatSeg{n+1}\) und damit \(n+1\) Glieder.

\FormulaThmDeltaK[Typisierung der Endpunkte eines Walks]{%
  \GraphStruct{V,E}\dsep\DirectedWalk{V}{E}{p}{n}\vdash
  p(0)\in V\land p(n)\in V%
}{DirectedWalkEndpointTyping}{%
  \DeltaRow{Mengen}{V\dsep E\dsep p\dsep n}%
}
\begin{tabproofwide}
  \proofstepwidestar[1]{\GraphStruct{V,E}}{\rA}
  \proofstepwidestar[2]{\DirectedWalk{V}{E}{p}{n}}{\rA}
  \proofstepwidestar[1,2]{%
    \begin{aligned}[t]
      &n\in\mathbb N\land p\colon\NatSeg{n+1}\to V\land{}\\[-2pt]
      &\forall i\in\NatSeg n\,((p(i),p(i+1))\in E)
    \end{aligned}}{\FormulaRefAuto{DirectedWalkCharacterization}[thm]{1,2}}
  \proofstepwidestar[1,2]{n\in\mathbb N}{\rAEa{3}}
  \proofstepwidestar[1,2]{p\colon\NatSeg{n+1}\to V}{%
    \FormulaRefAuto{P\land Q\land R\vdash Q}[thm]{3}}
  \proofstepwidestar[]{0\in\mathbb N}{\FormulaRefAuto{PeanoZero}[ax]}
  \proofstepwidestar[1,2]{0\leq n}{%
    \FormulaRefAuto{PeanoZeroLeq}[thm]{4}}
  \proofstepwidestar[1,2]{%
    0\leq n\leftrightarrow0\in\NatSeg{n+1}}{%
    \FormulaRefAuto{PeanoNatSegLeqAddOneMembership}[thm]{6,4}}
  \proofstepwidestar[1,2]{0\in\NatSeg{n+1}}{%
    \FormulaRefAuto{P\leftrightarrow Q, P\vdash Q}[thm]{8,7}}
  \proofstepwidestar[1,2]{n\leq n}{%
    \FormulaRefAuto{PeanoLeqReflexive}[thm]{4}}
  \proofstepwidestar[1,2]{%
    n\leq n\leftrightarrow n\in\NatSeg{n+1}}{%
    \FormulaRefAuto{PeanoNatSegLeqAddOneMembership}[thm]{4,4}}
  \proofstepwidestar[1,2]{n\in\NatSeg{n+1}}{%
    \FormulaRefAuto{P\leftrightarrow Q, P\vdash Q}[thm]{11,10}}
  \proofstepwidestar[1,2]{p(0)\in V}{%
    \FormulaRefAuto{F\colon A\to B\dsep x\in A\vdash F(x)\in B}[thm]{5,9}}
  \proofstepwidestar[1,2]{p(n)\in V}{%
    \FormulaRefAuto{F\colon A\to B\dsep x\in A\vdash F(x)\in B}[thm]{5,12}}
  \proofstepwidestar[1,2]{p(0)\in V\land p(n)\in V}{\rAI{13,14}}
\end{tabproofwide}

\FormulaThmDeltaK[Jeder Schritt eines Walks ist ein Bogen]{%
  \GraphStruct{V,E}\dsep\DirectedWalk{V}{E}{p}{n}\dsep
  i\in\NatSeg n\vdash(p(i),p(i+1))\in E%
}{DirectedWalkStepEdge}{%
  \DeltaRow{Mengen}{V\dsep E\dsep p\dsep n\dsep i}%
}
\begin{tabproofwide}
  \proofstepwidestar[1]{\GraphStruct{V,E}}{\rA}
  \proofstepwidestar[2]{\DirectedWalk{V}{E}{p}{n}}{\rA}
  \proofstepwidestar[3]{i\in\NatSeg n}{\rA}
  \proofstepwidestar[1,2]{%
    \begin{aligned}[t]
      &n\in\mathbb N\land p\colon\NatSeg{n+1}\to V\land{}\\[-2pt]
      &\forall j\in\NatSeg n\,((p(j),p(j+1))\in E)
    \end{aligned}}{\FormulaRefAuto{DirectedWalkCharacterization}[thm]{1,2}}
  \proofstepwidestar[1,2]{%
    \forall j\in\NatSeg n\,((p(j),p(j+1))\in E)}{%
    \FormulaRefAuto{P\land Q\land R\vdash R}[thm]{4}}
  \proofstepwidestar[1,2,3]{(p(i),p(i+1))\in E}{%
    \rRE{3,\rUE{5}}}
\end{tabproofwide}

\FormulaDefDeltaK[Endlicher gerichteter Pfad]{%
  \DirectedPath{V}{E}{p}{n}%
}{DirectedPathDef}{%
  \DeltaRow{Mengen}{V\dsep E\dsep p\dsep n}
  \DeltaPrem{Graph}{\GraphStruct{V,E}}
  \DeltaRow{\textbf{Axiome}}{}
  \DeltaRow{Walkstruktur}{%
    \DirectedWalk{V}{E}{p}{n}}
  \DeltaRow{Injektive Knotenfolge}{%
    p\colon\NatSeg{n+1}\inj V}
  \DeltaRow{\textbf{Neues Symbol}}{}
  \DeltaPrem{Struktur}{\DirectedPath{V}{E}{p}{n}}
}

\FormulaAxiomDeltaK[Walkstruktur]{%
  \begin{aligned}[t]
    &\GraphStruct{V,E}\dsep \DirectedPath{V}{E}{p}{n}\vdash{}\\[-2pt]
    &\DirectedWalk{V}{E}{p}{n}
  \end{aligned}%
}{DirectedPathWalkAxiom}{%
  \DeltaRow{Mengen}{V\dsep E\dsep p\dsep n}
}

\FormulaAxiomDeltaK[Injektive Knotenfolge]{%
  \begin{aligned}[t]
    &\GraphStruct{V,E}\dsep \DirectedPath{V}{E}{p}{n}\vdash{}\\[-2pt]
    &p\colon\NatSeg{n+1}\inj V
  \end{aligned}%
}{DirectedPathInjectiveAxiom}{%
  \DeltaRow{Mengen}{V\dsep E\dsep p\dsep n}
}

\FormulaThmDeltaK[Charakterisierung durch die Axiome]{%
  \GraphStruct{V,E}\vdash
  \DirectedPath{V}{E}{p}{n}\eqvdash
  \DirectedWalk{V}{E}{p}{n}\land
  p\colon\NatSeg{n+1}\inj V%
}{DirectedPathCharacterization}{%
  \DeltaRow{Mengen}{V\dsep E\dsep p\dsep n}
}
\begin{tabproofsplitwide}
  \proofpartwide{Hinrichtung}
    \proofstepwidestar[1]{\GraphStruct{V,E}}{\rA}
    \proofstepwidestar[2]{\DirectedPath{V}{E}{p}{n}}{\rA}
    \proofstepwidestar[1,2]{\DirectedWalk{V}{E}{p}{n}}{%
      \FormulaRefAuto{DirectedPathWalkAxiom}[ax]{1,2}}
    \proofstepwidestar[1,2]{p\colon\NatSeg{n+1}\inj V}{%
      \FormulaRefAuto{DirectedPathInjectiveAxiom}[ax]{1,2}}
    \proofstepwidestar[1,2]{%
      \begin{aligned}[t]
      &\DirectedWalk{V}{E}{p}{n}\land{}\\[-2pt]
      &p\colon\NatSeg{n+1}\inj V
    \end{aligned}}{\rAIStar{3,4}}
  \closeproofpartwide

  \proofpartwide{Rückrichtung}
    \proofstepwidestar[1]{\GraphStruct{V,E}}{\rA}
    \proofstepwidestar[2]{%
      \begin{aligned}[t]
      &\DirectedWalk{V}{E}{p}{n}\land{}\\[-2pt]
      &p\colon\NatSeg{n+1}\inj V
    \end{aligned}}{\rA}
    \proofstepwidestar[2]{\DirectedWalk{V}{E}{p}{n}}{\rAEa{2}}
    \proofstepwidestar[2]{p\colon\NatSeg{n+1}\inj V}{\rAEb{2}}
    \proofstepwidestar[1,2]{\DirectedPath{V}{E}{p}{n}}{%
      \FormulaRefAuto{DirectedPathDef}[def]{1,3,4}}
  \closeproofpartwide
\end{tabproofsplitwide}

\FormulaThmDeltaK[Funktionscharakter eines Pfads]{%
  \GraphStruct{V,E}\dsep\DirectedPath{V}{E}{p}{n}\vdash
  p\colon\NatSeg{n+1}\to V%
}{DirectedPathFunction}{%
  \DeltaRow{Mengen}{V\dsep E\dsep p\dsep n}
}
\begin{tabproof}
  \proofstep{1}{\GraphStruct{V,E}}{\rA}
  \proofstep{2}{\DirectedPath{V}{E}{p}{n}}{\rA}
  \proofstep{1,2}{p\colon\NatSeg{n+1}\inj V}{%
    \FormulaRefAuto{DirectedPathInjectiveAxiom}[ax]{1,2}}
  \proofstep{1,2}{p\colon\NatSeg{n+1}\to V}{%
    \FormulaRefAuto{Injektive Funktion}[def]{3}}
\end{tabproof}

\FormulaDefDeltaK[Gerichteter Kreis]{%
  \DirectedCycle{V}{E}{p}{n}%
}{DirectedCycleDef}{%
  \DeltaRow{Mengen}{V\dsep E\dsep p\dsep n\dsep i\dsep j}
  \DeltaPrem{Graph}{\GraphStruct{V,E}}
  \DeltaRow{\textbf{Axiome}}{}
  \DeltaRow{Natürliche Kantenlänge}{%
    n\in\mathbb N}
  \DeltaRow{Positive Kantenlänge}{%
    n\neq0}
  \DeltaRow{Walkstruktur}{%
    \DirectedWalk{V}{E}{p}{n}}
  \DeltaRow{Geschlossenheit}{%
    p(0)=p(n)}
  \DeltaRow{Verschiedene innere Knoten}{%
    \forall i,j\in\NatSeg n\,(p(i)=p(j)\rightarrow i=j)}
  \DeltaRow{\textbf{Neues Symbol}}{}
  \DeltaPrem{Struktur}{\DirectedCycle{V}{E}{p}{n}}
}

\FormulaAxiomDeltaK[Natürliche Kantenlänge]{%
  \begin{aligned}[t]
    &\GraphStruct{V,E}\dsep \DirectedCycle{V}{E}{p}{n}\vdash{}\\[-2pt]
    &n\in\mathbb N
  \end{aligned}%
}{DirectedCycleLengthAxiom}{%
  \DeltaRow{Mengen}{V\dsep E\dsep p\dsep n\dsep i\dsep j}
}

\FormulaAxiomDeltaK[Positive Kantenlänge]{%
  \begin{aligned}[t]
    &\GraphStruct{V,E}\dsep \DirectedCycle{V}{E}{p}{n}\vdash{}\\[-2pt]
    &n\neq0
  \end{aligned}%
}{DirectedCycleNonzeroAxiom}{%
  \DeltaRow{Mengen}{V\dsep E\dsep p\dsep n\dsep i\dsep j}
}

\FormulaAxiomDeltaK[Walkstruktur]{%
  \begin{aligned}[t]
    &\GraphStruct{V,E}\dsep \DirectedCycle{V}{E}{p}{n}\vdash{}\\[-2pt]
    &\DirectedWalk{V}{E}{p}{n}
  \end{aligned}%
}{DirectedCycleWalkAxiom}{%
  \DeltaRow{Mengen}{V\dsep E\dsep p\dsep n\dsep i\dsep j}
}

\FormulaAxiomDeltaK[Geschlossenheit]{%
  \begin{aligned}[t]
    &\GraphStruct{V,E}\dsep \DirectedCycle{V}{E}{p}{n}\vdash{}\\[-2pt]
    &p(0)=p(n)
  \end{aligned}%
}{DirectedCycleClosedAxiom}{%
  \DeltaRow{Mengen}{V\dsep E\dsep p\dsep n\dsep i\dsep j}
}

\FormulaAxiomDeltaK[Verschiedene innere Knoten]{%
  \begin{aligned}[t]
    &\GraphStruct{V,E}\dsep \DirectedCycle{V}{E}{p}{n}\vdash{}\\[-2pt]
    &\forall i,j\in\NatSeg n\,(p(i)=p(j)\rightarrow i=j)
  \end{aligned}%
}{DirectedCycleDistinctAxiom}{%
  \DeltaRow{Mengen}{V\dsep E\dsep p\dsep n\dsep i\dsep j}
}

\FormulaThmDeltaK[Charakterisierung durch die Axiome]{%
  \GraphStruct{V,E}\vdash
  \begin{aligned}[t]
    &\DirectedCycle{V}{E}{p}{n}\\[-2pt]
    &\quad\eqvdash n\in\mathbb N\land n\neq0\land
      \DirectedWalk{V}{E}{p}{n}\land p(0)=p(n)\land{}\\[-2pt]
    &\qquad\forall i,j\in\NatSeg n\,
        \bigl(p(i)=p(j)\rightarrow i=j\bigr)
  \end{aligned}%
}{DirectedCycleCharacterization}{%
  \DeltaRow{Mengen}{V\dsep E\dsep p\dsep n\dsep i\dsep j}
}
\begin{tabproofsplitwide}
  \proofpartwide{Hinrichtung}
    \proofstepwidestar[1]{\GraphStruct{V,E}}{\rA}
    \proofstepwidestar[2]{\DirectedCycle{V}{E}{p}{n}}{\rA}
    \proofstepwidestar[1,2]{n\in\mathbb N}{%
      \FormulaRefAuto{DirectedCycleLengthAxiom}[ax]{1,2}}
    \proofstepwidestar[1,2]{n\neq0}{%
      \FormulaRefAuto{DirectedCycleNonzeroAxiom}[ax]{1,2}}
    \proofstepwidestar[1,2]{\DirectedWalk{V}{E}{p}{n}}{%
      \FormulaRefAuto{DirectedCycleWalkAxiom}[ax]{1,2}}
    \proofstepwidestar[1,2]{p(0)=p(n)}{%
      \FormulaRefAuto{DirectedCycleClosedAxiom}[ax]{1,2}}
    \proofstepwidestar[1,2]{\forall i,j\in\NatSeg n\,(p(i)=p(j)\rightarrow i=j)}{%
      \FormulaRefAuto{DirectedCycleDistinctAxiom}[ax]{1,2}}
    \proofstepwidestar[1,2]{%
      \begin{aligned}[t]
      &n\in\mathbb N\land{}\\[-2pt]
      &n\neq0\land{}\\[-2pt]
      &\DirectedWalk{V}{E}{p}{n}\land{}\\[-2pt]
      &p(0)=p(n)\land{}\\[-2pt]
      &\forall i,j\in\NatSeg n\,(p(i)=p(j)\rightarrow i=j)
    \end{aligned}}{\rAIStar{3,4,5,6,7}}
  \closeproofpartwide

  \proofpartwide{Rückrichtung}
    \proofstepwidestar[1]{\GraphStruct{V,E}}{\rA}
    \proofstepwidestar[2]{%
      \begin{aligned}[t]
      &n\in\mathbb N\land{}\\[-2pt]
      &n\neq0\land{}\\[-2pt]
      &\DirectedWalk{V}{E}{p}{n}\land{}\\[-2pt]
      &p(0)=p(n)\land{}\\[-2pt]
      &\forall i,j\in\NatSeg n\,(p(i)=p(j)\rightarrow i=j)
    \end{aligned}}{\rA}
    \proofstepwidestar[2]{n\in\mathbb N}{\rAEa{2}}
    \proofstepwidestar[2]{n\neq0}{\rAEa{\rAEb{2}}}
    \proofstepwidestar[2]{\DirectedWalk{V}{E}{p}{n}}{\rAEa{\rAEb{\rAEb{2}}}}
    \proofstepwidestar[2]{p(0)=p(n)}{\rAEa{\rAEb{\rAEb{\rAEb{2}}}}}
    \proofstepwidestar[2]{\forall i,j\in\NatSeg n\,(p(i)=p(j)\rightarrow i=j)}{\rAEb{\rAEb{\rAEb{\rAEb{2}}}}}
    \proofstepwidestar[1,2]{\DirectedCycle{V}{E}{p}{n}}{%
      \FormulaRefAuto{DirectedCycleDef}[def]{1,3,4,5,6,7}}
  \closeproofpartwide
\end{tabproofsplitwide}

\FormulaDefDeltaK[Azyklischer gerichteter Graph]{%
  \AcyclicGraphStruct{V,E}%
}{DirectedAcyclicGraphDef}{%
  \DeltaRow{Mengen}{V\dsep E\dsep p\dsep n}
  \DeltaRow{\textbf{Axiome}}{}
  \DeltaRow{Graphstruktur}{%
    \GraphStruct{V,E}}
  \DeltaRow{Kreisfreiheit}{%
    \neg\exists n\in\mathbb N\,\exists p\,\DirectedCycle{V}{E}{p}{n}}
  \DeltaRow{\textbf{Neues Symbol}}{}
  \DeltaPrem{Struktur}{\AcyclicGraphStruct{V,E}}
}

\FormulaAxiomDeltaK[Graphstruktur]{%
  \begin{aligned}[t]
    &\AcyclicGraphStruct{V,E}\vdash{}\\[-2pt]
    &\GraphStruct{V,E}
  \end{aligned}%
}{DirectedAcyclicGraphGraphAxiom}{%
  \DeltaRow{Mengen}{V\dsep E\dsep p\dsep n}
}

\FormulaAxiomDeltaK[Kreisfreiheit]{%
  \begin{aligned}[t]
    &\AcyclicGraphStruct{V,E}\vdash{}\\[-2pt]
    &\neg\exists n\in\mathbb N\,\exists p\,\DirectedCycle{V}{E}{p}{n}
  \end{aligned}%
}{DirectedAcyclicGraphNoCycleAxiom}{%
  \DeltaRow{Mengen}{V\dsep E\dsep p\dsep n}
}

\FormulaThmDeltaK[Charakterisierung durch die Axiome]{%
  \AcyclicGraphStruct{V,E}\eqvdash
  \GraphStruct{V,E}\land
  \neg\exists n\in\mathbb N\,\exists p\,
    \DirectedCycle{V}{E}{p}{n}%
}{DirectedAcyclicGraphCharacterization}{%
  \DeltaRow{Mengen}{V\dsep E\dsep p\dsep n}
}
\begin{tabproofsplitwide}
  \proofpartwide{Hinrichtung}
    \proofstepwidestar[1]{\AcyclicGraphStruct{V,E}}{\rA}
    \proofstepwidestar[1]{\GraphStruct{V,E}}{%
      \FormulaRefAuto{DirectedAcyclicGraphGraphAxiom}[ax]{1}}
    \proofstepwidestar[1]{\neg\exists n\in\mathbb N\,\exists p\,\DirectedCycle{V}{E}{p}{n}}{%
      \FormulaRefAuto{DirectedAcyclicGraphNoCycleAxiom}[ax]{1}}
    \proofstepwidestar[1]{%
      \begin{aligned}[t]
      &\GraphStruct{V,E}\land{}\\[-2pt]
      &\neg\exists n\in\mathbb N\,\exists p\,\DirectedCycle{V}{E}{p}{n}
    \end{aligned}}{\rAIStar{2,3}}
  \closeproofpartwide

  \proofpartwide{Rückrichtung}
    \proofstepwidestar[1]{%
      \begin{aligned}[t]
      &\GraphStruct{V,E}\land{}\\[-2pt]
      &\neg\exists n\in\mathbb N\,\exists p\,\DirectedCycle{V}{E}{p}{n}
    \end{aligned}}{\rA}
    \proofstepwidestar[1]{\GraphStruct{V,E}}{\rAEa{1}}
    \proofstepwidestar[1]{\neg\exists n\in\mathbb N\,\exists p\,\DirectedCycle{V}{E}{p}{n}}{\rAEb{1}}
    \proofstepwidestar[1]{\AcyclicGraphStruct{V,E}}{%
      \FormulaRefAuto{DirectedAcyclicGraphDef}[def]{2,3}}
  \closeproofpartwide
\end{tabproofsplitwide}

\section{Erreichbarkeit}

\FormulaDefDeltaK[Gerichtete Erreichbarkeit]{%
  \GraphStruct{V,E}\vdash
  \begin{aligned}[t]
    u\mathrel{\GraphReach{E}}v\coloneqq{}&
      u\in V\land v\in V\land\Bigl(u=v\lor{}\\[-2pt]
      &\exists n\in\mathbb N\,\exists p\,
        \bigl(\DirectedWalk{V}{E}{p}{n}\land
          p(0)=u\land p(n)=v\bigr)\Bigr)
  \end{aligned}%
}{DirectedReachabilityDef}{%
  \DeltaRow{Mengen}{V\dsep E\dsep u\dsep v\dsep p\dsep n}
  \DeltaRow{\textbf{Neues Symbol}}{}
  \DeltaPrem{Erreichbarkeit}{\GraphReach{E}}
}

Die Gleichheitsalternative ist die explizite Länge-null-Konvention. Sie
macht jeden Knoten von sich selbst erreichbar, ohne für diesen Sonderfall
einen zusätzlichen Namen für die konstante Einpunktfolge einzuführen.

\FormulaThmDeltaK[Erreichbarkeit typisiert ihre Endpunkte]{%
  \GraphStruct{V,E}\dsep u\mathrel{\GraphReach{E}}v
  \vdash u\in V\land v\in V%
}{DirectedReachabilityEndpointTyping}{%
  \DeltaRow{Mengen}{V\dsep E\dsep u\dsep v}%
}
\begin{tabproof}
  \proofstep{1}{\GraphStruct{V,E}}{\rA}
  \proofstep{2}{u\mathrel{\GraphReach{E}}v}{\rA}
  \proofstep{1,2}{%
    \begin{aligned}[t]
      &u\in V\land v\in V\land\Bigl(u=v\lor{}\\[-2pt]
      &\exists n\in\mathbb N\,\exists p\,
        (\DirectedWalk{V}{E}{p}{n}\land p(0)=u\land p(n)=v)\Bigr)
    \end{aligned}}{\FormulaRefAuto{DirectedReachabilityDef}[def]{1,2}}
  \proofstep{1,2}{u\in V}{\rAEa{3}}
  \proofstep{1,2}{v\in V}{%
    \FormulaRefAuto{P\land Q\land R\vdash Q}[thm]{3}}
  \proofstep{1,2}{u\in V\land v\in V}{\rAI{4,5}}
\end{tabproof}

\FormulaThmDeltaK[Reflexivität der Erreichbarkeit auf dem Knotenträger]{%
  \GraphStruct{V,E}\dsep u\in V\vdash
  u\mathrel{\GraphReach{E}}u%
}{DirectedReachabilityReflexive}{%
  \DeltaRow{Mengen}{V\dsep E\dsep u}%
}
\begin{tabproofwide}
  \proofstepwidestar[1]{\GraphStruct{V,E}}{\rA}
  \proofstepwidestar[2]{u\in V}{\rA}
  \proofstepwidestar[]{u=u}{\rII}
  \proofstepwidestar[]{%
    u=u\lor\exists n\in\mathbb N\,\exists p\,
      (\DirectedWalk{V}{E}{p}{n}\land p(0)=u\land p(n)=u)}{\rOIa{3}}
  \proofstepwidestar[2]{%
    \begin{aligned}[t]
      &u\in V\land u\in V\land\Bigl(u=u\lor{}\\[-2pt]
      &\exists n\in\mathbb N\,\exists p\,
        (\DirectedWalk{V}{E}{p}{n}\land p(0)=u\land p(n)=u)\Bigr)
    \end{aligned}}{%
    \FormulaRefAuto{P\dsep Q\dsep R\vdash P\land Q\land R}[thm]{2,2,4}}
  \proofstepwidestar[1,2]{u\mathrel{\GraphReach{E}}u}{%
    \FormulaRefAuto{DirectedReachabilityDef}[def]{1,5}}
\end{tabproofwide}

\end{document}
